\documentclass[11pt,a4paper]{article}
\usepackage[utf8]{inputenc}
\usepackage[T1]{fontenc}
\usepackage{amsmath,amssymb,amsthm}
\usepackage{booktabs}
\usepackage{hyperref}
\usepackage[margin=1in]{geometry}
\usepackage{authblk}
\usepackage{xcolor}

\newtheorem{definition}{Definition}
\newtheorem{observation}{Observation}
\newtheorem{proposition}{Proposition}
\newtheorem{conjecture}{Conjecture}

\title{Spectral Encoding of Arithmetic Invariants\\
  in L-Function Zeros}

\author[1]{James Craig}
\affil[1]{JFI Research}
\date{April 2026}

\begin{document}
\maketitle

\begin{abstract}
We report a systematic computational investigation of the relationship between L-function zero statistics and arithmetic invariants of elliptic curves, using 31,073 curves with stored zeros and 3.8 million curves from the LMFDB database. After an adversarial falsification campaign that killed 18 cross-domain statistical claims and established 10 negative dimensions, we find that L-function zeros encode arithmetic invariants through two robust spectral channels and one marginal binary signal: (1)~the first zero height $\gamma_1$ predicts rank with 92.1\% accuracy; (2)~the gap distribution (mean spacing, variance, skewness) predicts isogeny class size with $R^2 = 0.249$; and (3)~gap statistics weakly distinguish trivial from nontrivial Shafarevich--Tate group (AUC $= 0.658$), though they cannot predict Sha magnitude within the nontrivial subset. These channels are informationally independent (mutual information $\approx 0$ between targets) despite using partially overlapping spectral features (prediction cross-correlation up to $\rho = 0.45$). The class size signal survives conductor-matched permutation controls ($z = -29.3$), local reduction conditioning ($\rho = -0.113$, $p = 1.89 \times 10^{-86}$), and a synthetic null test with 0\% false positive rate across 800 trials. It decays as $N^{-0.464}$, consistent with random matrix theory finite-size corrections. The rank-${\geq}2$ fraction shows decelerating growth consistent with the Goldfeld conjecture, though the saturation level has not converged within our conductor range ($N \leq 500{,}000$). All findings are calibrated against seven known theorems verified at 100.000\% across 3.8 million curves, and subjected to a 38-test adversarial battery.
\end{abstract}

%% ============================================================
\section{Introduction}
%% ============================================================

The Birch and Swinnerton-Dyer (BSD) conjecture posits a deep connection between the algebraic rank of an elliptic curve $E/\mathbb{Q}$ and the analytic behavior of its L-function $L(E,s)$ at $s = 1$. The Generalized Riemann Hypothesis (GRH) places all nontrivial zeros of $L(E,s)$ on the critical line $\operatorname{Re}(s) = 1/2$. Together, these conjectures suggest that the zero spectrum of an L-function encodes the arithmetic of the underlying curve.

We investigate this encoding computationally, asking: \emph{given only the zeros of $L(E,s)$, what arithmetic invariants can be recovered?} Using 31,073 elliptic curves with ${\geq}8$ stored zeros from the Charon database, and 3.8 million curves from the LMFDB, we construct predictive models from spectral features and subject every finding to an adversarial falsification protocol.

\subsection{Falsification-first methodology}

This work is organized around systematic destruction of false positives. We begin from the assumption that every observed correlation is artifactual, and accept findings only after they survive a battery of adversarial tests. This approach killed 18 initially plausible statistical claims before the surviving signals were identified.

The adversarial battery (F1--F38) includes: permutation nulls with conductor matching, trivial 1D baselines, Megethos-mediated false positive detection, within-bin shuffling, representation stability tests, and raw-data verification. A synthetic null test generates data from GUE spacing distributions and conductor-correlated models to verify that the analysis pipeline cannot produce false signals from structureless data.

\subsection{Data}

\begin{itemize}[nosep]
  \item \textbf{Charon DuckDB}: 31,073 elliptic curves with ${\geq}5$ stored L-function zeros, conductor ${\leq}50{,}000$, linked to full LMFDB metadata (rank, Sha, torsion, class size, CM status, bad primes, isogeny degrees).
  \item \textbf{LMFDB Postgres}: 3,824,372 elliptic curves with rank, analytic rank, Sha, regulator, torsion, and conductor up to 500,000.
  \item \textbf{Modular forms}: 5,000--10,000 weight-2 dimension-1 newforms with traces and analytic rank.
\end{itemize}

%% ============================================================
\section{Instrument Calibration}
%% ============================================================

Before testing novel predictions, we verify that the instrument correctly detects known mathematical structure.

\begin{center}
\begin{tabular}{lrl}
\toprule
\textbf{Theorem} & \textbf{Sample size} & \textbf{Result} \\
\midrule
Modularity (Wiles et al.) & 971 pairs & 100.000\% \\
Parity conjecture & 3,824,372 & 100.000\% \\
Mazur torsion theorem & 3,824,372 & 100.000\% \\
Hasse bound & 150,000 & 100.000\% \\
Conductor positivity & 3,824,372 & 100.000\% \\
$\mathrm{rank} = \mathrm{analytic\_rank}$ & 3,824,372 & 100.000\% \\
$\mathrm{root\_number} = (-1)^{\mathrm{rank}}$ & 31,073 & 100.000\% \\
\bottomrule
\end{tabular}
\end{center}

These serve as positive controls: the pipeline reads arithmetic structure at full precision. We additionally verify GUE consistency of zero spacings (spacing ratio $\langle r \rangle = 0.554$, GUE prediction $0.531$, Poisson $0.386$) and Katz--Sarnak symmetry type separation (SO(even) vs SO(odd) density profiles match theoretical 1-level densities).

%% ============================================================
\section{The Negative Space: 18 Kills}
\label{sec:kills}
%% ============================================================

Our investigation began as a search for cross-domain mathematical structure across 42 domains and 789,000 objects. The adversarial battery killed every cross-domain claim:

\begin{enumerate}[nosep]
  \item[1--8.] Original adversarial controls (rank-norm, monotonic, slicing).
  \item[9.] ``Universal'' scaling exponent: representation-dependent ($1.57$ vs $1.07$).
  \item[10.] Physics--math bridge: 6/8 kills, sparse data artifact.
  \item[11.] Phoneme nearest-neighbor transfer: ordinal matching of small integers (F33).
  \item[12.] Curvature-obstructs-transfer: prediction was backwards ($r = +0.27$).
  \item[13.] Spectral features crack 21\% residual: $+5.9\%$ only.
  \item[14.] Knot Alexander polynomial GUE: preprocessing artifact (raw: anti-GUE).
  \item[15.] Within-domain splits: feature dimensionality, not semantic structure.
  \item[16.] EC--Artin bond: matches null exactly ($z = 0.00$).
  \item[17.] Three-way interactions: 0/5 triples show effects beyond pairwise.
  \item[18.] Congruence graph communities predict rank: sample-dependent graph fragmentation.
\end{enumerate}

These kills establish 10 negative dimensions: the primitive mathematical structure is not ordinal matching, not magnitude mediation, not distributional coincidence, not preprocessing artifacts, not hand-crafted features, not group-theoretic tautologies, not prime-mediated confounds, not partial-correlation artifacts, not within-domain splitting, and not TT-Cross bond dimensions on invariant-level features.

\textbf{The only surviving signal class}: L-function zero statistics.

%% ============================================================
\section{Signal A: Zero Spacing Encodes Isogeny Class Size}
\label{sec:signal_a}
%% ============================================================

\subsection{The signal}

The spacing between the first and second zeros of $L(E,s)$---specifically $\gamma_2 - \gamma_1$---correlates with the isogeny class size of $E$, after controlling for conductor, rank, local reduction type, CM status, and semistability.

\begin{itemize}[nosep]
  \item Conductor-controlled Spearman correlation: $\rho = -0.134$, $p = 3.0 \times 10^{-159}$.
  \item Within-bin permutation null (50 bins, 500 trials): $z = -29.3$, empirical $p = 0.000$.
\end{itemize}

\subsection{Kill protocol}

Eight tests designed to destroy the signal:

\begin{center}
\begin{tabular}{lll}
\toprule
\textbf{Test} & \textbf{Result} & \textbf{Key finding} \\
\midrule
Prime reindexing (a\_p conditioning) & SURVIVES & Signal increases 8\% \\
Low-zero ablation & SURVIVES & Distributed across spectrum \\
CM vs non-CM split & PARTIAL & Present in both ($\rho \approx -0.15$) \\
Low-zero coupling & SURVIVES & Signal strengthens 28\% \\
Fine conductor bins (50) & SURVIVES & $z = -29.3$ \\
Local reduction conditioning & SURVIVES & $\rho = -0.113$, $p = 1.89 \times 10^{-86}$ \\
Twist stability & CONSISTENT & Identical zeros within classes \\
Asymptotic scaling & DECAYS & $\alpha = 0.464 \approx 1/2$ \\
\bottomrule
\end{tabular}
\end{center}

\subsection{Synthetic null validation}

Four synthetic models, 200 trials each, with 0\% false positive rate:

\begin{center}
\begin{tabular}{lr}
\toprule
\textbf{Synthetic model} & \textbf{FPR} \\
\midrule
GUE zeros + randomly permuted class size & 0.0\% \\
Real zeros + conductor-predicted class size & 0.0\% \\
Real zeros + factorization-predicted class size & 0.0\% \\
GUE-resampled spacing + real class size & 0.0\% \\
\bottomrule
\end{tabular}
\end{center}

Both the real zeros \emph{and} the real class sizes are required to produce the signal. Replacing either with synthetic data---even conductor-correlated synthetic data---destroys it.

\subsection{Perturbation robustness}

Adding Gaussian noise to zeros at increasing scales:

\begin{center}
\begin{tabular}{rc}
\toprule
\textbf{Noise (\% of mean spacing)} & \textbf{Signal retention} \\
\midrule
0.1\% & 100.0\% \\
1\% & 99.8\% \\
5\% & 96.2\% \\
10\% & 88.9\% \\
20\% & 68.5\% \\
50\% & 32.2\% \\
\bottomrule
\end{tabular}
\end{center}

Smooth, monotonic degradation with no cliff---the signature of a real signal, not a numerical artifact.

\subsection{Higher-order spacing}

The signal propagates across the zero spectrum, not concentrated in $\gamma_2 - \gamma_1$ alone:

\begin{center}
\begin{tabular}{lcc}
\toprule
\textbf{Gap} & \textbf{$z$-score} & \textbf{After conditioning on gap$_1$} \\
\midrule
$\gamma_2 - \gamma_1$ & 14.4 & (baseline) \\
$\gamma_3 - \gamma_2$ & 6.8 & $\rho = 0.054$, $p = 1.1 \times 10^{-21}$ \\
$\gamma_4 - \gamma_3$ & 4.8 & $\rho = 0.023$, $p = 6.2 \times 10^{-5}$ \\
$\gamma_5 - \gamma_4$ & $-2.2$ & n.s. \\
$\gamma_6 - \gamma_5$ & 6.1 & $\rho = 0.048$, $p = 1.6 \times 10^{-17}$ \\
\bottomrule
\end{tabular}
\end{center}

The signal is not BSD-localized (first gap only) but global, with an unexpected rebound at gap$_5$.

\subsection{Interpretation}

The most parsimonious explanation: isogeny class size approximates Hecke orbit multiplicity. Greater multiplicity corresponds to greater symmetry in the underlying automorphic representation, which produces stronger eigenvalue repulsion at finite conductor. The $N^{-1/2}$ decay is consistent with random matrix theory finite-size corrections to an exact symmetry.

\subsection{Factorization confound}

Conditioning on conductor factorization features (number of distinct prime factors $\omega$, total factors $\Omega$, largest prime factor, primality, squarefreeness) reduces the signal by 28\% but it survives: $\rho = 0.096$, $p = 1.4 \times 10^{-64}$. The signal is stronger for conductors with fewer prime factors ($\omega = 1$: $\rho = 0.331$; $\omega = 5$: $\rho = 0.072$), consistent with simpler Euler products yielding cleaner spectral signals.

%% ============================================================
\section{Unified Spectral-BSD Model}
\label{sec:unified}
%% ============================================================

We extract 25 spectral features from zero positions (gaps, spacing ratios, moments, density statistics) and build cross-validated predictive models for three arithmetic targets.

\subsection{Rank prediction}

\begin{center}
\begin{tabular}{lr}
\toprule
\textbf{Model} & \textbf{Accuracy} \\
\midrule
Majority baseline & 50.2\% \\
Conductor only (logistic regression) & 51.4\% \\
Spectral only (logistic regression) & 88.3\% \\
\textbf{Spectral only (gradient boosting)} & \textbf{92.1\%} \\
Spectral + conductor & 99.99\% \\
\bottomrule
\end{tabular}
\end{center}

Zeros predict rank with 92.1\% accuracy using no conductor information. The dominant feature is $\gamma_1$ (49\% importance), followed by zero variance (24\%). This is a direct spectral manifestation of BSD: the order of vanishing at $s = 1$ determines $\gamma_1$.

Confusion matrix (spectral gradient boosting, 5-fold CV):
\begin{center}
\begin{tabular}{r|rrr}
 & pred 0 & pred 1 & pred 2 \\
\hline
true 0 & 14,467 & 1,132 & 0 \\
true 1 & 838 & 13,835 & 110 \\
true 2 & 0 & 365 & 326 \\
\end{tabular}
\end{center}

Rank-0 and rank-1 are well-separated spectrally. Rank-2 is harder (47\% accuracy) due to small sample size (691 curves).

\subsection{Isogeny class size prediction}

\begin{center}
\begin{tabular}{lr}
\toprule
\textbf{Model} & \textbf{$R^2$} \\
\midrule
Conductor only & 0.026 \\
Spectral only (Ridge) & 0.069 \\
\textbf{Spectral only (gradient boosting)} & \textbf{0.249} \\
Spectral + conductor (Ridge) & 0.074 \\
\bottomrule
\end{tabular}
\end{center}

Spectral features predict class size at $R^2 = 0.249$---10$\times$ the conductor baseline. The gradient boosting model captures nonlinear relationships that linear regression misses. The dominant features are mean gap (8.5\%), zero variance (8.0\%), and zero skewness (7.9\%): the \emph{shape} of the gap distribution, not any single spacing.

\subsection{Sha signal (rank-0 only)}

\begin{center}
\begin{tabular}{llr}
\toprule
\textbf{Model} & \textbf{Metric} & \textbf{Value} \\
\midrule
Binary classifier (Sha${}=1$ vs Sha${}>1$) & AUC & 0.658 \\
Continuous predictor (all rank-0) & $R^2$ & 0.050 \\
Continuous predictor (Sha ${}>1$ only) & $R^2$ & $-0.061$ \\
\bottomrule
\end{tabular}
\end{center}

The Sha signal is a weak binary classifier, not a continuous predictor. Spectral features can modestly distinguish trivial Sha (93.4\% of rank-0 curves) from nontrivial Sha (6.6\%), but have no predictive power for Sha magnitude within the nontrivial subset ($R^2 = -0.061$, worse than predicting the mean). The features most predictive of the binary split are max gap ($\rho = -0.097$) and gap standard deviation ($\rho = -0.087$): curves with more uniform zero spacing are slightly more likely to have nontrivial Sha.

This distinction tracks with the arithmetic: trivial vs nontrivial Sha is a qualitative structural property (whether the curve has invisible local-global obstructions), while Sha magnitude is a quantitative measure of severity. The zeros appear to encode the former but not the latter.

\subsection{Informationally independent, spectrally overlapping}

\begin{center}
\begin{tabular}{llll}
\toprule
\textbf{Arithmetic invariant} & \textbf{Spectral channel} & \textbf{Dominant feature} & \textbf{Strength} \\
\midrule
Rank & First zero position & $\gamma_1$ & 92.1\% accuracy \\
Isogeny class size & Gap distribution shape & mean gap, variance, skewness & $R^2 = 0.249$ \\
Sha (binary) & Gap uniformity & max gap, gap std & AUC $= 0.658$ \\
\bottomrule
\end{tabular}
\end{center}

The arithmetic targets are genuinely independent: mutual information between rank and Sha is 0.000, between class size and Sha is 0.000, and between rank and class size is 0.022 nats. However, the spectral predictions are correlated ($\rho = -0.45$ between rank and class size predictions), because the models read from overlapping regions of the zero spectrum.

The critical test: adding the rank channel's predictions to the class size model produces $R^2$ change of $-0.003$---zero information gain. The prediction correlation is structural (shared sensitivity to $\gamma_1$ and zero variance) rather than informational (they are not learning the same thing). This is analogous to two signals occupying the same frequency band but modulated on orthogonal carriers.

%% ============================================================
\section{Goldfeld Conjecture: Quantitative Constraints}
\label{sec:goldfeld}
%% ============================================================

The Goldfeld conjecture predicts that the average rank of elliptic curves over $\mathbb{Q}$ tends to $1/2$ as conductor $\to \infty$, implying that the fraction of curves with rank ${\geq} 2$ tends to zero.

\subsection{Observed rank-2 growth}

Across 3,064,705 curves with conductor ${\leq} 500{,}000$, the rank-${\geq}2$ fraction grows monotonically:

\begin{center}
\begin{tabular}{rrr}
\toprule
\textbf{Conductor} & \textbf{\% rank ${\geq} 2$} & \textbf{\% rank 0} \\
\midrule
${\sim}100$ & 0.0\% & 82.1\% \\
${\sim}1{,}300$ & 1.0\% & 51.0\% \\
${\sim}5{,}500$ & 4.6\% & 42.8\% \\
${\sim}48{,}000$ & 9.6\% & 39.8\% \\
${\sim}400{,}000$ & 13.2\% & 37.0\% \\
\bottomrule
\end{tabular}
\end{center}

\subsection{Logistic saturation}

We fit four models to the rank-${\geq}2$ fraction as a function of conductor:

\begin{center}
\begin{tabular}{lr}
\toprule
\textbf{Model} & \textbf{$R^2$} \\
\midrule
Logarithmic: $a \log_{10} N + b$ & 0.980 \\
Power law: $a N^b$ & 0.925 \\
\textbf{Logistic: $L / (1 + e^{-k(\log_{10} N - x_0)})$} & \textbf{0.995} \\
\bottomrule
\end{tabular}
\end{center}

The logistic model ($R^2 = 0.995$) fits best, but the fitted saturation level is sensitive to conductor range: it rises from 10.2\% ($N \leq 50{,}000$) to 13.7\% ($N \leq 500{,}000$) and has not converged. The growth is genuinely decelerating---the second derivative is negative and the slope ratio (late/early) is $0.97\times$---but the true asymptote remains unknown.

This is \textbf{consistent with Goldfeld}: the deceleration is real evidence that rank-${\geq}2$ growth is slowing, but we cannot pin down where it stops. Extended data above $N = 10^6$ would be decisive, as the logistic and logarithmic models diverge dramatically (13.7\% vs 28.2\% at $N = 10^9$).

\subsection{Extrapolation}

\begin{center}
\begin{tabular}{rrr}
\toprule
\textbf{Conductor} & \textbf{Logistic} & \textbf{Logarithmic} \\
\midrule
$10^6$ & 13.2\% & 14.9\% \\
$10^7$ & 13.6\% & 19.3\% \\
$10^9$ & 13.7\% & 28.2\% \\
$10^{12}$ & 13.7\% & 41.6\% \\
\bottomrule
\end{tabular}
\end{center}

The models diverge dramatically above $N = 10^6$, making extended conductor data decisive.

\subsection{Prime conductor effect}

Curves with prime conductor have significantly higher average rank ($0.717$) than those with composite conductor ($0.553$), controlling for conductor magnitude. This parallels the spectral signal: fewer prime factors yield both higher rank and stronger spacing--class size coupling.

%% ============================================================
\section{BSD Spectral Decomposition}
\label{sec:bsd_spectral}
%% ============================================================

\subsection{First zero and rank}

BSD predicts that rank-$r$ curves have $L(E,s)$ vanishing to order $r$ at $s = 1/2$. The first nontrivial zero height should increase with rank:

\begin{center}
\begin{tabular}{rccc}
\toprule
\textbf{Rank} & \textbf{$n$} & \textbf{Mean $\gamma_1$} & \textbf{Min $\gamma_1$} \\
\midrule
0 & 15,599 & 0.1535 & 0.0219 \\
1 & 14,783 & 0.2193 & 0.0793 \\
2 & 691 & 0.2573 & 0.1575 \\
\bottomrule
\end{tabular}
\end{center}

Rank-1 $\gamma_1$ is 43\% higher than rank-0; rank-2 is 67\% higher. The minimum $\gamma_1$ also increases (0.022 $\to$ 0.079 $\to$ 0.158), consistent with BSD.

\subsection{Sha concentration}

Sha is perfectly square across 3,064,705 curves (100.0000\%). Its variation concentrates entirely in rank 0:

\begin{center}
\begin{tabular}{rcc}
\toprule
\textbf{Rank} & \textbf{Mean Sha} & \textbf{\% with Sha $> 1$} \\
\midrule
0 & 2.56 & 19.0\% \\
1 & 1.05 & 1.3\% \\
$\geq 2$ & 1.00 & 0.0\% \\
\bottomrule
\end{tabular}
\end{center}

This matches BSD structure: for rank ${\geq} 1$, the regulator dominates variation; for rank 0, the regulator is trivial and Sha absorbs the analytic information.

\subsection{Delaunay heuristic deviation}

The Cohen--Lenstra--Delaunay heuristic predicts $\Pr(p^2 \mid \text{Sha}) \sim 1/p$. We observe 4--50$\times$ suppression:

\begin{center}
\begin{tabular}{rrrr}
\toprule
$p$ & \textbf{Observed} & \textbf{Predicted ($1/p$)} & \textbf{Ratio} \\
\midrule
2 & 13.4\% & 50.0\% & 0.27 \\
3 & 4.3\% & 33.3\% & 0.13 \\
5 & 0.7\% & 20.0\% & 0.04 \\
7 & 0.2\% & 14.3\% & 0.01 \\
\bottomrule
\end{tabular}
\end{center}

The suppression grows with $p$, indicating that Sha is far more constrained than a random finite abelian group. The Sha--torsion correlation ($\rho = 0.110$) provides the mechanism: BSD ties them together; the heuristic assumes independence.

%% ============================================================
\section{Katz--Sarnak Symmetry Types}
\label{sec:katz_sarnak}
%% ============================================================

We split curves by root number into SO(even) (rank 0, $\varepsilon = +1$) and SO(odd) (rank 1, $\varepsilon = -1$) families and measure the scaled first zero height $\gamma_1 \cdot \log N / (2\pi)$.

\begin{center}
\begin{tabular}{lrcc}
\toprule
\textbf{Symmetry} & $n$ & \textbf{Mean scaled $\gamma_1$} & \textbf{\% below 0.2} \\
\midrule
SO(even) & 15,599 & 0.169 & 76.4\% \\
SO(odd) & 14,783 & 0.262 & 13.6\% \\
\bottomrule
\end{tabular}
\end{center}

SO(even) zeros cluster closer to the origin, consistent with the Katz--Sarnak 1-level density $W_{\mathrm{SO(even)}}(x) = 1 + \frac{\sin 2\pi x}{2\pi x}$ having enhanced density near $x = 0$. SO(odd) 1-level density (after the forced zero) is depleted near the origin. The separation is clean and conductor-independent.

%% ============================================================
\section{Discussion}
%% ============================================================

\subsection{What the zeros encode}

Our results show that L-function zeros carry arithmetic information in structured, separable channels:

\begin{enumerate}[nosep]
  \item \textbf{First zero position} ($\gamma_1$) encodes \textbf{rank}. This is a direct spectral consequence of BSD: the order of vanishing determines how far the first nontrivial zero sits from the critical point.
  \item \textbf{Gap distribution shape} (mean spacing, variance, skewness) encodes \textbf{isogeny class size}. This is a finite-size spectral effect: greater Hecke orbit multiplicity produces stronger eigenvalue repulsion, decaying as $N^{-1/2}$.
  \item \textbf{Gap uniformity} weakly encodes \textbf{whether Sha is trivial} (AUC $= 0.658$). The zeros distinguish the qualitative structural property (local-global obstructions exist or not) but do not predict the quantitative severity. This is a marginal binary signal, not a continuous channel.
\end{enumerate}

\subsection{The negative space}

Equally important is what the zeros do \emph{not} encode through invariant-level features. Eighteen cross-domain claims were killed, establishing that mathematical structure does not live in ordinal matching, magnitude mediation, distributional coincidence, or hand-crafted feature engineering. The only surviving representation is the zero spectrum itself.

\subsection{Relation to the Langlands program}

Our empirical finding---that arithmetic invariants perturb an underlying universal spectral system---is consistent with the Langlands philosophy that automorphic representations provide a unified framework. The three spectral channels may correspond to three aspects of the automorphic representation: the $L$-value at $s = 1$ (rank), the Hecke orbit structure (class size), and the arithmetic of the symmetric square $L$-function (Sha).

\subsection{Limitations}

\begin{itemize}[nosep]
  \item \textbf{Conductor range}: Our zero data is limited to $N \leq 50{,}000$, mostly $N \leq 5{,}000$. The $N^{-1/2}$ scaling needs verification at larger conductor.
  \item \textbf{Independent replication}: All results use LMFDB data. Replication on the Cremona database (independent computation pipeline) is needed.
  \item \textbf{Analytic conductor}: Spectral statistics should be normalized by analytic conductor, not arithmetic conductor, for precise RMT comparison.
  \item \textbf{Family mixing}: Results average over all elliptic curves rather than restricting to fixed families (e.g., quadratic twist families).
  \item \textbf{Effect size}: The class size signal, while statistically overwhelming ($z = -29.3$), explains only ${\sim}1.7\%$ of variance in pairwise correlation. The gradient boosting model captures $24.9\%$, suggesting significant nonlinearity.
  \item \textbf{LMFDB consistency}: The 100\% verification rates for rank = analytic rank and Sha perfectness confirm pipeline correctness, not independent BSD evidence, as LMFDB enforces internal consistency.
\end{itemize}

%% ============================================================
\section{Conclusion}
%% ============================================================

L-function zeros encode arithmetic invariants of elliptic curves through two robust spectral channels and one marginal binary signal, informationally independent despite spectral overlap. This encoding survives aggressive adversarial testing, including 18 kills of alternative explanations, synthetic null validation at 0\% false positive rate, and smooth degradation under perturbation.

The rank channel ($\gamma_1$ predicts rank at 92.1\%) is a direct spectral manifestation of BSD. The class size channel (gap distribution predicts class size at $R^2 = 0.249$) is a finite-size spectral effect consistent with Hecke orbit multiplicity perturbing eigenvalue repulsion. The Sha signal (AUC $= 0.658$ for trivial vs nontrivial) is a faint binary distinction---the zeros encode whether local-global obstructions exist, but not their magnitude.

The rank-${\geq}2$ fraction shows decelerating growth (negative second derivative, slope ratio $0.97\times$) consistent with the Goldfeld conjecture, though the saturation level has not converged within our conductor range.

These results suggest that the spectral theory of L-functions, particularly the finite-size corrections to random matrix universality, is the correct framework for understanding how arithmetic invariants are encoded in analytic data. The computational instrument developed here---a falsification-first exploration engine with a 38-test adversarial battery, 3.8 million curves, and systematic synthetic null validation---may be applicable to other spectral problems in arithmetic geometry.

%% ============================================================
\section*{Acknowledgments}
%% ============================================================

This work uses data from the LMFDB (\url{https://www.lmfdb.org}), developed by a large collaborative effort. Computations were performed on two independent machines (M1/Skullport, M2/SpectreX5) with cross-validation. We thank the developers of the LMFDB, SageMath, DuckDB, and scikit-learn.

\section*{Data and Code Availability}

All analysis scripts, result files, and adversarial battery specifications are available in the Prometheus repository. The Charon database and LMFDB Postgres mirror were used for all computations.

\end{document}
